\chapter{C1.1 with $k^2\eta$ a unit: a complete contact certificate}
\label{component:p95350f69a353}
\noindent\textbf{Component }\path{p95350f69a353}\par
{\small\noindent Source: \path{R13/advances/c1_small/q_unit_completion.tex}\par}
{\small\noindent SHA256: \path{95350f69a35347503869de108b3124c5ad72b0daadde632ce44dcbb244e890df}\par}
{\small\noindent Inventory use: SC10--SC20 and complete full-normal/contact proof | C1\_\allowbreak{}BARE\_\allowbreak{}OTHER\_\allowbreak{}TERMINALS; precise subordinate uses retained below\par}
\medskip
\section{C1.1 with $k^2\eta$ a unit: a complete contact certificate}
\label{p95350f69a353:sec:r7-c1-q-unit}
\begin{theorem}
On the original C1.1 source, suppose $q=k^2\eta$ is a positive bounded
unit. The entire original parent is tensor SOS at every independent
projective test and every real or complex state. All positive opening
rates and the permitted literal boundary $D=0$ are included.
\end{theorem}
The proof uses the exact full quadratic (SC3), not its leading scalar.
The $q$-infinitesimal part of C1.1 is not included in this theorem.

\paragraph{One more centered packet.}
Use the variables of (SC1)--(SC5), and put $s_D=J/R$. Retain the
additional nonnegative square
\[
                  4D_b(J_b+s_DH)^2.                 \tag{SC10}
\]
It vanishes at actual contact, where $H=RvZ$ and $J_b=-JvZ$.
Its removal replaces $4D_bJ_b^2$ in (SC3) exactly by
\[
                  -8D_bs_DHJ_b-4D_bs_D^2H^2.
\]
In particular it removes the entire standalone $(vZ)^2$ term.
This step allows $D_b=0$ literally and does not divide by $D_b$.

Set $Z=1$ on the remaining state chart. Put $c=\sqrt q$ and
$w=(W+uY)/R$, and use the four separated positive monomials
\[
 A=\frac{cuw}{b^2},\quad B=\frac{cvw}{bg},\quad
 C=\frac{uv}{b},\quad D=c^2g,\quad
 \Gamma=(ABCD)^{1/4}.
 \tag{SC11}
\]
Their exact contact is $(u,v,w)=(k,cg,bg)$, with $k=bc$, and
$A=B=C=D$ there. Notice that $H=RA/c$, $L=RB/c$, $V=C$,
$g=D/q$, and $Hv=R\Gamma^2/q$.

Here are explicit formulas for the entire scalar at $N=N_y=0$,
after (SC10). Write $a_t=1-\rho_0+pv^2$. Define
\begin{align*}
 P_H={}&4D_v\{c_b(1-\rho_0)^2+a_b\rho_0^2/D_u\}
       +2b^2\{a_t^2+v^2\rho_0^2/D_u\}\\
 &+4\alpha_bg^2b^4v^2\rho_0^2/D_u
       -8D_bs_Db^2a_t-4D_bs_D^2,\\
 P_L={}&4[h_2(1-p)+h_bpD_v]/R,\\
 P_{Hv}={}&-4[D_u(1+pv^2)-u^2\rho_0]
             -8\alpha_bg^2b^2\rho_0+8D_bs_DD_u.
\end{align*}
The exact scalar is
\[
 f=f_AA^2+f_BB^2+f_CC^2+f_DD^2-f_\Gamma\Gamma^2,
 \tag{SC12}
\]
where
\[
 f_A=R^2P_H/q,\quad f_B=R^2P_L/q,\quad
 f_C=2D_u,\quad f_D=4\alpha_bD_u/q^2,
 \quad f_\Gamma=-RP_{Hv}/q.
\]
These are rational functions of $u,v$ over bounded source coefficients,
with only positive-unit denominators. For example,
\[
 c_b(1-\rho_0)^2=
   \frac{K(1+b^2r)E(1-r)}{2P_dR^2},\qquad
 a_b\rho_0^2=\frac{P_4K^2r}{2P_dR^2}.
\]
Thus neither $r$ opening survives in their denominators.
For $u,v,b,g$ infinitesimal their residues are
\[
               \overline f_A=\overline f_B=
 \overline f_C=\overline f_D=2,\qquad
               \overline f_\Gamma=8.               \tag{SC13}
\]
This uses the exact source residues in (SC2), including
$\overline D_b=(1-2\overline q)/2$. It does not assume a positive
lower bound on $D_b$.

Here is an explicit finite denominator certificate used below.
Every function just displayed is $P(u,v)/(1+u^2)$ for a polynomial
$P$ with bounded pure-source coefficient array. For a monomial use
\[
 x^m-x_0^m=(x-x_0)\sum_{j=0}^{m-1}x^{m-1-j}x_0^j.
\]
Applying this identity to each monomial supplies finite first and
second difference arrays with bounded source coefficients. For the
denominator use only the exact identity
\[
 \frac1{1+u^2}-\frac1{1+k^2}
 =-\frac{(u-k)(u+k)}{(1+u^2)(1+k^2)}.
\]
Thus the first and second finite Taylor arrays about $(k,cg)$
are bounded and use only pure-testing $1+u^2$ and pure-source
$1+k^2$ denominators, both units. The same construction after
$u=U^2,v=V^2$ gives the half-variable arrays below. This is a
finite polynomial identity construction, not an analytic series.

\paragraph{Every first normal and the complete hard block.}
Let $\delta_\rho=\rho_0-\rho$. The full first-normal coefficients
of the remaining quadratic, with no normal suppressed, are
\begin{align*}
 \mathcal L_N={}&\frac8R(h_bD_v-h_2)L
       +4buvg(a_t-2D_bs_D)H-4D_uvgV,\\
 \mathcal L_y={}&8D_v\sqrt{s}\,\delta_\rho H\\
 &+\frac H{\sqrt s}\left[
 -\frac{8D_va_bu^2\rho_0}{D_u}+8D_bs_Db^2
 +\frac{8\alpha_bg^2b^4v^2\rho_0}{D_u}
 -4b^2a_t+\frac{4b^2v^2\rho_0}{D_u}\right]\\
 &+\frac{4bu}{\sqrt{s}}V
       -\frac{8\alpha_bgb^2v}{\sqrt{s}}(gZ).
                                                        \tag{SC14}
\end{align*}
Each is a bounded infinitesimal-coefficient linear combination of
$A,B,C,D$. At actual contact it is zero exactly, since the source
is a stationary zero and every removed packet vanishes there.
The coefficients of $N^2,N_y^2,NN_y$ are respectively
\[
\begin{aligned}
 H_{NN}&=4(h_2+h_bD_v)/R+2g^2u^2v^2,\\
 H_{yy}&=4D_v\left(1-\frac{a_bu^2}{D_us}\right)
       +\frac{4\alpha_bg^2b^4v^2}{D_us}
       +\frac{2b^2}{s}\left(1+\frac{v^2}{D_u}\right),\\
 H_{Ny}&=-4buvg/\sqrt{s}.
\end{aligned}                                                     \tag{SC15}
\]
Thus the hard Gram has residue $\operatorname{diag}(4q,4)$.
In (SC15), $H_{Ny}$ denotes the polynomial coefficient of $NN_y$,
not the off-diagonal matrix entry.

\paragraph{A finite geometric-mean square lift.}
Here is a detailed lift for the particular raw module (SC12).
It is not a claim about arbitrary contact-vanishing polynomials.
Define $\alpha_i=A_i-\Gamma$, with
$(A_1,A_2,A_3,A_4)=(A,B,C,D)$, and put
\[
 T_4=\sum_iA_i^2-4\Gamma^2.
\]
A fixed positive Gram on the following finite roots represents $T_4$:
\[
 \alpha_i,\qquad
 Q_{ij|kl}=\sqrt{A_iA_j}-\sqrt{A_kA_l},\qquad
 \sqrt\Gamma(\sqrt{A_i}-\sqrt{A_j}),\qquad
 \sqrt\Gamma((A_iA_j)^{1/4}-(A_kA_l)^{1/4}).          \tag{SC16}
\]
Use all four $\alpha_i$, all three pair partitions, and all six
unordered differences in the third group. To verify strictness, average
positive multiples of the symmetric circuit decomposition in (SC9) and
\[
 T_4=\sum_i\alpha_i^2+2\Gamma\left(\sum_iA_i-4\Gamma\right).
\]

over the three decompositions
\[
 \sum_iA_i-4\Gamma=(\sqrt{A_i}-\sqrt{A_j})^2
 +(\sqrt{A_k}-\sqrt{A_l})^2
 +2((A_iA_j)^{1/4}-(A_kA_l)^{1/4})^2.
\]
The additional differences $A_i-A_j$ in the first decomposition are
$\alpha_i-\alpha_j$ and give nonnegative corrections to the already
strict $\alpha$ block.

Write $\epsilon_i=f_i-2$, $\epsilon_\Gamma=f_\Gamma-8$, and
$\sigma=\sum_i\epsilon_i-\epsilon_\Gamma$. Then
\[
 f-2T_4=\sum_i\epsilon_i\alpha_i^2
       +2\Gamma\sum_i\epsilon_i\alpha_i+\Gamma^2\sigma.
 \tag{SC17}
\]
All $\epsilon$ have zero residue. The exact contact value implies
$\sigma(k,cg)=0$. Clear the pure-testing denominator $D_u$ if needed;
monic division gives
$\sigma=(u-k)\sigma_u+(v-cg)\sigma_v$ with bounded coefficients.
For the monomials (SC11), the exact conductors are
\[
\begin{aligned}
 \Gamma(u-k)&=\sqrt{uk}\,Q_{13|24},\\
 \Gamma(v-cg)&=\sqrt{vcg}\,Q_{23|14},\\
 \Gamma(w-bg)&=\sqrt{wbg}\,Q_{12|34}.
\end{aligned}                                                     \tag{SC18}
\]
Thus the last term in (SC17) becomes $\Gamma$ times a linear
combination of the $Q$ roots with infinitesimal bounded coefficients.
The exact identity
\[
 Q_{ij|kl}-\tfrac12(A_i+A_j-A_k-A_l)
 =-\tfrac12(\sqrt{A_i}-\sqrt{A_j})^2
  +\tfrac12(\sqrt{A_k}-\sqrt{A_l})^2               \tag{SC19}
\]
changes these terms into infinitesimal quadratic corrections on (SC16)
and $\Gamma\sum_i\mu_i\alpha_i$, where each $\mu_i$ is infinitesimal
and is rational in $\sqrt u,\sqrt v$ with bounded coefficient arrays
and positive-unit denominators.

The exact three contact jets imply that the four source values
$\mu_i(k,cg)$ coincide. Indeed the differential vectors of
$\alpha_i/\Gamma$ at contact are the four rows
\[
 (1/2,-1/2,1/2),\quad(-1/2,1/2,1/2),\quad
 (1/2,1/2,-1/2),\quad(-1/2,-1/2,-1/2)
\]
in logarithmic coordinates $(u/k,v/(cg),w/(bg))$. Their sole linear
relation is the sum of all four rows. Cancelling the nonzero source
factors in this equality establishes coincidence; no inverse of a
small number enters a Gram coefficient.
Call the common infinitesimal value $\mu_*$. The common term is
\[
 \mu_*\Gamma\sum_i\alpha_i
   =\mu_*\Gamma\left(\sum_iA_i-4\Gamma\right),
\]
an infinitesimal correction on the explicitly displayed variance roots.
For each remaining $\mu_i-\mu_*$, monic division in $\sqrt u,\sqrt v$
has bounded coquotients. Apply the half-variable conductors
\[
\begin{aligned}
 \Gamma(\sqrt u-\sqrt k)
   &=(uk)^{1/4}\sqrt\Gamma
          ((AC)^{1/4}-(BD)^{1/4}),\\
 \Gamma(\sqrt v-\sqrt{cg})
   &=(vcg)^{1/4}\sqrt\Gamma
          ((BC)^{1/4}-(AD)^{1/4}).
\end{aligned}                                                     \tag{SC20}
\]
Their coefficient factors are infinitesimal. Multiplying by the
remaining $\alpha_i$ proves that the entire scalar correction in
(SC17), not merely its low-order jet, is a bounded zero-residue
quadratic form on (SC16). This uses only two finite monic divisions.
The boundedness of their coefficient arrays follows from the explicit
rational functions (SC12), whose test denominators are powers of
$D_u$, and from the displayed square-root substitutions; bounded
function values alone are not being used as a substitute for this fact.

The same argument for a first normal is shorter. If
$\mathcal L=\sum_i a_iA_i$, write
$\mathcal L=\sum_i a_i\alpha_i+\Gamma\sum_i a_i$.
Its exact contact zero gives $\sum_i a_i(k,cg)=0$; one monic division
and (SC18) allocate it on (SC16), with infinitesimal coefficients.
Apply this to both complete expressions (SC14).
Together with (SC15), this yields a bounded strictly positive-residue
Gram for the entire original remainder. Restore (SC10), the paired
packets, and all positive pure-source and pure-testing normalizations.

\paragraph{Exhaustion and legal radicals.}
For $q$ a unit, $e^2=q/g^2$ is unbounded. The exact adjacent-event
rotations and the $W$-unit state argument (SC8) therefore reduce every
unadmitted negative-core test to $u,v,W$ infinitesimal and $Z$ a unit.
Literal axes are direct original-core exits. Use exact sign gauges and
normalize $Z=1$. The inherited pure-testing Schur identity forces
$W+uY>0$ on the remaining negative-core chart: with
$w_t=(W+uY)/D_u$, the eliminated scalar is
\[
 \frac{uvD_uD_v}{\mathcal D}
       (D_uuv w_t^2-2D_uD_vw_t+uv),\qquad
 \mathcal D=u^2v^4+u^2v^2+u^2+v^4+2v^2+1>0,
\]
which is nonnegative when $w_t\le0$; its eliminated two-square
matrix is strictly positive. Hence $w=(W+uY)/R>0$.
Every monomial in (SC11), and every product under a radical in
(SC16)--(SC20), factors into a positive pure-source quantity and a
positive pure-testing quantity. The occurrence of $R$ in $w$ is a
positive pure-source unit, not a mixed denominator. All test
coefficients remain bounded because $u,v,w$ are infinitesimal.
If the general toric contact lemma is used, its third coordinate
is literally the pure-testing $W+uY$, with pure-source contact
$Rbg$; this direction is inactive in every coefficient function.
Complex state evaluations split into real and imaginary parts.
This supplies every independent test and proves the theorem.

\paragraph{Audit scope.}
The parent and the scoped toric-lift reviewer independently reviewed
the finite lift, bounded rational arrays, full normal block, legal
radicals, and joins within the same agent system. The thirty-nine
exact checks verify the displayed whole-polynomial and source
identities. A further six exact polynomial checks verify the actual
source contact value, all three scalar first jets, and both full
normal contact zeros directly from (NC1)--(NC4). Their explicit
denominator-cleared numerators are printed in
\path{check_q_unit_contact.py}; they confirm the original
stationary-zero premise without sampling. No numerical table is a
premise of this argument.
